
% SGA 3，第 VII 篇_B；英文重建本第一轮的中文译本。
% 范围： 节 3 标题、开头约定及 节 3.1 全文,
% 依据文献局部页 66--67, 止于下节之前： 节 3.2.
% 合并阅读本页： 579--580. 重版印刷页： 568--569.
% 范围内原版边码： 531.
% 依据： Polo--Gille Exp7B-13oct24.pdf,
% SHA-256 D303690E0195654BBE7C3DAA5D89EA23E24CD709A3575217234AC19AE8836694.
% 参照了 Jacob Reinhold 的 CC BY 4.0 Markdown，作为对照来源。
% 本组件无图；所有数学内容均为原生 LaTeX。
\section*{3. 特征 \(0\) 特有的现象}
\phantomsection
\label{sga3:VIIB:3}
\addcontentsline{toc}{section}{3. 特征 0 特有的现象}
在\SGARefLink{sga3:VIIB:3}{第 3 节}中，始终假设伪紧环 \(k\) 包含有理数域 \(\mathbb Q\)。
\subsection*{3.1. 引理}
\phantomsection
\label{sga3:VIIB:3.1}
\noindent\textbf{引理 3.1。}
\phantomsection
\label{sga3:VIIB:3.1:lemma}
{\itshape 设 \(C\) 是交换含幺 \(\mathbb Q\)-代数，\(L\) 是 \(C\) 上的 Lie 代数，其底 \(C\)-模为自由模。则典范映射
\[
  L\longrightarrow U(L)
\]
将 \(L\) 与 \(U(L)\) 的本原元素集合等同。}
\textit{证明。}
将 \(L\) 与其在 \(U(L)\) 中的典范像等同。设 \(I\) 是全序集，且
\[
  (x_i)_{i\in I}
\]
为 \(L\) 的一组以 \(I\) 为指标集的基。用 \(\mathbb N^{(I)}\) 表示非负整数族
\[
  n=(n_i)_{i\in I}
\]
的集合，其中除可能有限多个指标
\[
  i_1<i_2<\cdots<i_s;
\]
以外都有 \(n_i=0\)；这些指标依赖于 \(n\)。置
\[
  x^n
  =
  x_{i_1}^{\,n_{i_1}}
  x_{i_2}^{\,n_{i_2}}
  \cdots
  x_{i_s}^{\,n_{i_s}},
  \qquad
  n!
  =
  (n_{i_1}!)(n_{i_2}!)\cdots(n_{i_s}!).
\]
由 Poincaré--Birkhoff--Witt 定理，元素 \(x^n\) 构成 \(U(L)\) 的一组基。此外，
\begin{equation}
  \Delta_{U(L)}
  \left(\frac{x^n}{n!}\right)
  =
  \sum_{0\leq m\leq n}
  \frac{x^m}{m!}
  \otimes
  \frac{x^{\,n-m}}{(n-m)!}.
  \tag{*}
  \label{sga3:VIIB:3.1:pbw-comultiplication}
\end{equation}
这里求和遍历所有满足 \(0\leq m\leq n\) 的 \(m\in\mathbb N^{(I)}\)，即
\[
  0\leq m_i\leq n_i
  \qquad\text{对每个 }i.
\]
立即可得
\[
  \Delta_{U(L)}(u)=u\otimes1+1\otimes u
\]
当且仅当 \(u\) 是各 \(x_i\) 的线性组合。引理得证。
